% Imporditud: tools/import_eksperimendid.py
% Allikas: Eesti füüsikaolümpiaadi eksperimendi ülesannete
%   kogu (Rael Kalda, 2023), ülesanne Ü25, lahendus L25.
\setAuthor{EFO žürii}
\setRound{lõppvoor}
\setYear{2016}
\setNumber{P E2}
\setDifficulty{1}
\setTopic{vedelike mehaanika}
\setVahendid{pliiats, mõõtesilinder (50 ml), joogitops veega}
\setPunktid{12}
\setAllikas{Eksperimendi kogu Ü25, lahendus lk 43}
\setLahendusseis{toores}

\prob{Pliiats}
Määrake pliiatsi tihedus.

\hint

\solu
Asetame pliiatsi mõõtesilindrisse nii, et ta ujuks. Fikseerime ruumala muutuse $\Delta$$V_0$. Pliiatsi mass võrdub väljatõrjutud vedeliku massiga, seega m = $\rho$v$\Delta$$V_0$. Pliiatsi ruumala leidmisel paneme tähele, et meil vett ei ole piisavalt, et pliiatsit tervenisti uputada. Küll aga me saame teda uputada kahest erinevast otsast, märkides pliiatsil koha, milleni uputame. Mõõtesilindri abil mõõdame $\Delta$$V_1$ ja $\Delta$$V_2$. Pliiatsi ruumala seega V = $\Delta$$V_1$ + $\Delta$$V_2$. Tiheduseks saame

$\rho$ = $\rho$v $\Delta$$V_0$ . $\Delta$$V_1$ + $\Delta$$V_2$
\probend
